  \section{Conclusion}
  \label{sec:conclusion}

  This project added a natural-language advisory layer to an existing DT, the incubator
  of Gomes et al.\ \cite{gomes2025incubator}, without changing how the twin computes. The layer is
  two skills over tools that drive the twin's own models, and it answers questions about opening
  the lid and heating the box first. The research question asked what to do when the twin is asked
  something no tool answers directly, such as how long the lid may stay open: code the search as a
  new tool, or delegate it to the LLM over the existing tools.

  For this twin the answer is to delegate first and code what recurs. With skills that describe the
  kind of question, the model reached the reference on every inverse question, as a coded search
  did, again in 53 of the 54 repeats and in every run of two questions written after
  the skills, and invented no value where none could be given.
  Delegating took time of the same order but cost 194 to 473 thousand input tokens per answer
  (US\$0.25 to 0.42 at list prices), and up to US\$0.59 on a search over two inputs. It also gave
  up the guard at the simulation's resolution that the coded search applies by rule, and labelled
  its search late. With skills that described no
  search, the model asked back, declined or mostly stopped short of the limit. Without skills it
  searched at once and found most limits, but went past the twin, naming a room
  temperature computed on states the twin never published. A search is therefore worth coding when its question
  recurs, must clear the resolution or must give the same answer by construction, and is otherwise
  better left to the model.

  The project makes four contributions:

  \begin{itemize}
    \item An advisory layer for an existing twin made only of skills and wrapper tools, with the
      twin's code unchanged, and extended once, by the heater boost, without changing the first
      advisory (\Cref{sec:system}).
    \item An evaluation that scores a language layer against references computed by the same tools
      it calls, with a raw and a guarded answer at the simulation's resolution, and compares
      its cost with a coded search (\Cref{sec:method}).
    \item A measured comparison of delegating and coding the same searches, in closeness, sources,
      tokens, time and what had to be foreseen (\Cref{sec:results}).
    \item A rule for when to code and when to delegate, four ways in which the skills limit or free
      the model, and five steps for applying the approach to another twin
      (\Cref{sec:discussion}).
  \end{itemize}

  Four directions follow from the results. The split proposed in \Cref{sec:build-or-ask}, where a
  coded method searches a question the model formulates and the twin judges, could keep both the
  model's flexibility and the guarantees of code, and should be built and compared with the two
  ways tested here. Questions delegated often could then be promoted to coded tools, so that what
  the twin can answer grows from the questions it is asked. The lid-opening tool should report its
  resolution and its skill ask for the guard, to test whether one stated line closes the gap at
  the resolution. And the comparison needs wider evidence: other twins and models, more unseen
  questions and searches over several inputs, several scored runs per question and a second rater.
  Two further directions concern the twin rather than the layer: checking the open-lid disturbance
  against the recorded openings, and modelling the batch's own heat, which would change the
  references but not the method.

  An LLM can let a twin answer questions no one coded for, provided its skills allow the
  search and leave the judgement to the twin.
